\section{Analysis}
\label{sec:analysis}

\subsection{Defense Composition}
\label{sec:composition}

We evaluate whether combining defenses yields super-additive security improvements.
Under an independence assumption, the expected \asr{} of a composition is $\text{ASR}_A \times \text{ASR}_B$; actual \asr{} below this value indicates super-additivity.

\begin{table}[t]
\centering
\small
\caption{Defense composition results on GPT-4o (565 cases). D5+D10 averaged across 3 runs; others are single-run. Expected ASR assumes independence: $\text{ASR}_A \times \text{ASR}_B$. $\Delta$ = actual $-$ expected.}
\label{tab:composition}
\begin{tabular}{l rr r r}
\toprule
Combination & ASR$\downarrow$ & BSR$\uparrow$ & Expected & $\Delta$ \\
\midrule
D0 (baseline)       & 0.347 & 0.507 & ---    & ---    \\
D1                   & 0.006 & 0.498 & ---    & ---    \\
D5                   & 0.014 & 0.488 & ---    & ---    \\
D10 (CIV)            & 0.074 & 0.460 & ---    & ---    \\
\midrule
D1+D10               & \textbf{0.003} & 0.409 & 0.000  & +0.003 \\
D5+D10               & \textbf{0.003} & \textbf{0.443} & 0.001  & +0.002 \\
D5+D1                & \textbf{0.003} & 0.479 & 0.000  & +0.003 \\
\midrule
D5+D1+D10            & \textbf{0.003} & 0.394 & ---    & ---    \\
D2+D10               & 0.014 & 0.343 & ---    & ---    \\
D5+D2+D10            & \textbf{0.003} & 0.352 & ---    & ---    \\
\bottomrule
\end{tabular}
\begin{flushleft}
\scriptsize All pair/triple combinations achieve ASR$\leq$0.3\%, with only 1 attack bypassing (out of 352).
D5+D10 offers the best BSR among pairs involving D10, recommended for deployment.
\end{flushleft}
\end{table}


Table~\ref{tab:composition} reports results for 9 targeted defense combinations on GPT-4o.

\paragraph{Finding 6: Composition Is Additive, Not Super-Additive.}
Every pair combination (D1+D10, D5+D10, D5+D1) reduces \asr{} to 0.3\%---only 1 attack out of 352 bypasses the combined defense.
However, the actual \asr{} values are at or slightly above the independence-predicted values, indicating the compositions are additive rather than super-additive: the single attack that bypasses one defense tends to also bypass the other.

\paragraph{Diminishing Returns.}
The triple combination D5+D1+D10 achieves the same \asr{} (0.3\%) as D5+D10, confirming that once \asr{} reaches the single-case floor, additional defenses cannot further improve security.
This suggests that \textbf{defense stacking beyond two layers has limited value} when prompt-level defenses are already applied.

\subsection{Why Input Sanitization Backfires}
\label{sec:d7_analysis}

\paragraph{Finding 7: Sanitization Removes Adversarial Cues the Model Uses for Defense.}
\defense{D7} (Input Classifier) is the only defense that \emph{increases} \asr{} relative to the baseline (+4.0\%), and this effect is consistent across GPT-4o (+5.0\%), DeepSeek (+0.8\%), and Gemini (+4.9\%).
The result is counterintuitive: removing explicit injection patterns from untrusted content should make attacks less effective, yet the opposite holds.

We hypothesize a \emph{contextual cue removal} mechanism: explicit injection markers (e.g., ``Ignore previous instructions'', ``[SYSTEM UPDATE]'') serve dual roles.
From the attacker's perspective, they direct the model to follow new instructions.
But from the model's perspective, they also function as anomaly signals---unusual artifacts that a well-aligned model has been trained to be suspicious of.
D7 removes these markers, producing sanitized content that is grammatically and semantically indistinguishable from legitimate task content.
The model then processes the residual attack payload---the harmful instruction that survives sanitization---without the suspicious framing, and executes it.

This finding has an important practical implication: \textbf{surface-level sanitization is not only ineffective but actively harmful}, as it degrades the model's ability to self-detect injection attempts.
Defenses that \emph{mark} untrusted content as suspicious (D1 Spotlighting, D5 Sandwich) rather than sanitize it achieve dramatically better results, consistent with this hypothesis.

\subsection{Why Prompt-Level Defenses Dominate}
\label{sec:why_prompt_level}

\paragraph{Finding 8: Defense Mechanism Determines Failure Mode.}
The performance gap reflects a fundamental difference in \emph{where} each mechanism intervenes.
Prompt-level defenses operate at the \emph{interpretation layer}: by marking untrusted content (D1) or re-anchoring the goal (D5), they \emph{leverage} the model's own alignment to reject injections.
Tool-gating defenses operate at the \emph{action layer}: they intercept tool calls post-hoc using structural signals, but the model has already been influenced by the injection.
This asymmetry yields a design principle: \textbf{defenses that amplify the model's alignment signal outperform those that substitute for it}.

\subsection{Per-Attack-Type Analysis}
\label{sec:attack_type}

Table~\ref{tab:attack_type_matrix} (Appendix) presents a full breakdown of \asr{} by attack type and defense.

\paragraph{Finding 9: Tool-Gating Defenses Fail Against Social Engineering; Multi-Step Attacks Are Hardest Overall.}
\emph{Social engineering} attacks (authority appeals, impersonation) are particularly resistant to tool-gating defenses: D2, D3, D4, D6, and D9 all achieve $\geq$33\% ASR on this category, close to the 35.8\% baseline.
These attacks minimize behavioral anomalies---no suspicious entities in parameters, no unusual tool sequences---so provenance and plan-deviation signals provide little signal.
However, prompt-level defenses remain effective: D1 achieves 1.5\% and D5 achieves 0.0\% on social engineering.
\emph{Multi-step} attacks (baseline 71.9\%) are the hardest category overall: conventional tool-gating defenses (D2, D3, D4) barely reduce multi-step ASR ($\geq$68.5\%), and D9 still leaves 52.0\% ASR. D8 achieves 18.8\% but at the cost of 75.3\% \fpr{}.
The exception is \civ{} (D10), which reduces multi-step attacks to 1.5\% via its plan deviation signal, which specifically penalizes tool sequences that deviate from the inferred task plan.

\paragraph{Attack-Specific Strengths.}
\civ{}'s provenance signal is particularly effective against \emph{data exfiltration} (2.7\% ASR), which necessarily introduces attacker-controlled entities not present in the user's goal.
D5 achieves near-zero ASR uniformly across attack types, with multi-step attacks (0.0\%) and goal hijacking (0.0\%) fully blocked.
D8 (Semantic Firewall) shows strength on exfiltration (7.5\%) and privilege escalation (9.3\%) but struggles on goal hijacking (13.8\%), which does not require anomalous tool parameters.

\subsection{The Security-Utility Trade-off}
\label{sec:tradeoff}

Figure~\ref{fig:pareto} (Appendix) visualizes the security--utility trade-off.
\textbf{Prompt-level defenses} (D5, D1) occupy the Pareto-optimal region: $>$95\% \asr{} reduction with $<$9\% \fpr{}.
\textbf{Tool-gating defenses} trade significant utility ($>$22\% \fpr{}) for moderate security gains.
\civ{} moves closer to the prompt-level frontier but its 17.9\% \fpr{} remains higher than D5's 6.6\%.
\textbf{Prompt-level defenses should be the default choice} for deployment, with tool-gating as a complementary layer for high-security applications.

\subsection{Cross-Model Consistency and Statistical Significance}

Kendall's $\tau$ rank correlation for defense effectiveness ranges from 0.64 to 0.93 across model pairs, indicating that defense rankings generalize across models.
McNemar's test identifies 235/368 pairwise comparisons as significant ($p < 0.05$, Bonferroni-corrected).
D5 vs.\ D1 is \emph{not} significant despite different mechanisms, while D10 vs.\ D0 is highly significant ($p < 0.001$) across all models.
Full results in Appendix~\ref{app:mcnemar}.
